\subsection{Brecha uniforme del operador canónico reducido}\label{sec:P03}
Se usan exclusivamente las cotas deterministas de la Sección~\ref{sec:P02}.
Sea $n$ suficientemente grande para que
\[
 c/n\le\pi_{ij}\le M/n\le1/8,\qquad h_D,h_A\ge\eta>0.
\]
Los símbolos $c,M,\eta$ son constantes, no índices de casilla.
La matriz de incidencia $A$ tiene una fila por casilla y una columna por línea.
Ponemos
\[
 g_{ij}=\frac{\pi_{ij}(1-\pi_{ij})}{1-4\pi_{ij}},\quad
 S=\operatorname{diag}(0_R,0_C,h_D,h_A)+A^T\operatorname{diag}(g)A.
                                                        \tag{G1}
\]
Entonces $c/n\le g_{ij}\le2M/n$.
Se usan normas euclídeas ordinarias en las $6n-2$ líneas.

\begin{proposicion}
El núcleo de $S$ es exactamente
$\operatorname{span}\{v_0\}$, con
$v_0=(\mathbf1_R,-\mathbf1_C,0_D,0_A)$.
En $v_0^\perp$ se tiene
\[
 \gamma\|v\|^2\le v^TSv\le(1+8M)\|v\|^2,
 \qquad \gamma=\min\{c/2,\eta/5\}>0.               \tag{G2}
\]
Las constantes son independientes de $n$.
\end{proposicion}
\begin{proof}
Escribimos $v=(r_i,t_j,d_k,a_l)$ y
$s_{ij}=r_i+t_j+d_{i-j}+a_{i+j-1}$.
La condición $v\perp v_0$ equivale a $\sum_i r_i=\sum_jt_j$.
Por tanto
\[
 \frac1n\sum_{ij}(r_i+t_j)^2
 =\sum_i r_i^2+\sum_jt_j^2+
       \frac2n\left(\sum_i r_i\right)^2
 \ge\|r\|^2+\|t\|^2.
\]
La desigualdad $(s-d-a)^2\le2s^2+4(d^2+a^2)$
y el hecho de que cada diagonal tiene a lo sumo $n$ casillas
dan
\[
 \|v\|^2\le\frac2n\sum_{ij}s_{ij}^2
                        +5(\|d\|^2+\|a\|^2).     \tag{G3}
\]
Por (G1),
\[
 v^TSv\ge\frac cn\sum s_{ij}^2+
                          \eta(\|d\|^2+\|a\|^2),
\]
lo que prueba la cota inferior. Además
$s_{ij}^2\le4(r_i^2+t_j^2+d_{i-j}^2+a_{i+j-1}^2)$.
Cada línea tiene a lo sumo $n$ casillas, así que
$\sum s_{ij}^2\le4n\|v\|^2$. Como $h_D,h_A\le1$,
se obtiene la cota superior.
Finalmente $Sv_0=0$. Si $v^TSv=0$, la positividad de
$g,h_D,h_A$ obliga a $d=a=0$ y $r_i+t_j=0$ para todo $i,j$.
Luego $r$ es constante y $t$ su opuesto: no hay otros ceros.
\end{proof}

\paragraph{Consecuencia.}
La inversa de $S$ sobre $v_0^\perp$ tiene norma espectral
a lo sumo $\gamma^{-1}$, con $\gamma$ independiente de $n$.
